% ============================================================
%  Hiero Shared SDK Automations — Architecture Plan
%  Five-page architecture document | arc42-compliant
%  Draft for Sophie/explorer3 review | May 2026
% ============================================================
\documentclass[10pt, a4paper]{article}

% ── Packages ────────────────────────────────────────────────
\usepackage[T1]{fontenc}
\usepackage[utf8]{inputenc}
\usepackage[margin=2.3cm]{geometry}
\usepackage{setspace}
\usepackage{titlesec}
\usepackage{titling}
\usepackage{booktabs}
\usepackage{tabularx}
\usepackage{array}
\usepackage{longtable}
\usepackage{graphicx}
\usepackage{xcolor}
\usepackage{hyperref}
\usepackage{enumitem}
\usepackage{parskip}
\usepackage{fancyhdr}
\usepackage{mdframed}
\usepackage{etoolbox}

% ── Colour palette ──────────────────────────────────────────
\definecolor{hieroBlue}{HTML}{1A3A5C}
\definecolor{hieroAccent}{HTML}{2E75B6}
\definecolor{hieroLight}{HTML}{EAF2FB}
\definecolor{hieroGray}{HTML}{5A5A5A}
\definecolor{hieroRule}{HTML}{CCCCCC}
\definecolor{noteYellow}{HTML}{FFF8DC}
\definecolor{noteYellowBorder}{HTML}{C8A000}

% ── Hyperref setup ──────────────────────────────────────────
\hypersetup{
  colorlinks  = true,
  linkcolor   = hieroAccent,
  urlcolor    = hieroAccent,
  citecolor   = hieroAccent,
  pdftitle    = {Hiero Shared SDK Automations — Architecture Plan},
  pdfauthor   = {Draft for Sophie/explorer3},
}

% ── Section title styling ───────────────────────────────────
\titleformat{\section}
  {\large\bfseries\color{hieroBlue}}
  {\thesection}{1em}{}[\vspace{-2pt}\color{hieroRule}\rule{\linewidth}{0.4pt}]

\titleformat{\subsection}
  {\normalsize\bfseries\color{hieroBlue}}
  {\thesubsection}{1em}{}

\titleformat{\subsubsection}
  {\normalsize\bfseries\color{hieroGray}}
  {\thesubsubsection}{1em}{}

\titlespacing*{\section}{0pt}{13pt}{6pt}
\titlespacing*{\subsection}{0pt}{7pt}{3pt}

% ── Header / footer ─────────────────────────────────────────
\pagestyle{fancy}
\fancyhf{}
\renewcommand{\headrulewidth}{0.4pt}
\renewcommand{\footrulewidth}{0.4pt}
\fancyhead[L]{\small\color{hieroGray}\textit{Hiero Shared SDK Automations — Architecture Plan}}
\fancyhead[R]{\small\color{hieroGray}May 2026}
\fancyfoot[C]{\small\color{hieroGray}\thepage}

% ── Table helpers ───────────────────────────────────────────
\newcolumntype{L}[1]{>{\raggedright\arraybackslash}p{#1}}
\newcolumntype{R}[1]{>{\raggedleft\arraybackslash}p{#1}}

% ── Note box ────────────────────────────────────────────────
\newmdenv[
  backgroundcolor = noteYellow,
  linecolor       = noteYellowBorder,
  linewidth       = 0.8pt,
  leftmargin      = 0pt,
  rightmargin     = 0pt,
  innerleftmargin = 10pt,
  innerrightmargin= 10pt,
  innertopmargin  = 8pt,
  innerbottommargin=8pt,
  skipabove       = 8pt,
  skipbelow       = 8pt,
]{notebox}

% ── List style ──────────────────────────────────────────────
\setlist[itemize]{leftmargin=1.5em, itemsep=1.2pt, topsep=2.5pt}
\setlist[enumerate]{leftmargin=1.5em, itemsep=1.2pt, topsep=2.5pt}

% ── Compact tables ──────────────────────────────────────────
\AtBeginEnvironment{tabularx}{\small}
\renewcommand{\arraystretch}{0.92}

% ── No paragraph indent ─────────────────────────────────────
\setlength{\parindent}{0pt}
\setlength{\parskip}{3pt}

% ============================================================
\begin{document}
% ============================================================

% ── Title block ─────────────────────────────────────────────
\begin{center}
  {\LARGE\bfseries\color{hieroBlue} Hiero Shared SDK Automations}\\[6pt]
  {\large\color{hieroGray} Five-Page Architecture Plan}\\[4pt]
  {\small\color{hieroGray}
    arc42-compliant \quad|\quad
    Microkernel pattern \quad|\quad
    Draft for Sophie\,/\,explorer3 review \quad|\quad
    May 2026
  }
\end{center}

\vspace{4pt}
\color{hieroRule}\hrule\color{black}
\vspace{6pt}

\begin{notebox}
  \textbf{Core principle:} centralise reusable automation decision logic, not repository control.

  \textbf{Architecture pattern:} Microkernel Architecture (Richards, 2015). The core automation
  kernel resides in \texttt{packages/core}; GitHub Actions and Probot are delivery adapters.

  \textbf{arc42 coverage:} Sections 1-12 are addressed in this five-page plan. C4 views are
  included for L1 (Context), L2 (Container), and L3 (Component).
\end{notebox}

This document is intentionally architecture-first and follows the arc42 structure. C4 L1 is
shown as a diagram; L2 and L3 are text views. A separate migration strategy follows from this
architecture and is intentionally out of scope here.

% ============================================================
\section{Introduction and Goals}
% ============================================================

Hiero Shared SDK Automations provides a shared automation decision core while preserving SDK
repository autonomy over workflow triggers, permissions, and security controls. The audience is
SDK maintainers, automation contributors, and reviewers who need a concise, reviewable baseline.

\subsection{Top goals (quality drivers)}

\begin{itemize}[itemsep=0.6pt, topsep=2pt]
  \item Security: shared logic must not hide permissions, tokens, or runner hardening.
  \item Maintainability: fix automation logic once and adopt deliberately across SDKs.
  \item Evolvability: enable Actions now and Probot later without duplicating behavior.
  \item Testability: prove parity before write mode through fixtures and comparisons.
  \item Adoptability: each SDK adopts incrementally and can roll back by SHA.
\end{itemize}

\subsection{Stakeholders and expectations}

\begin{tabularx}{\linewidth}{@{} L{4cm} X @{}}
  \toprule
  \textbf{Stakeholder} & \textbf{Expectation} \\
  \midrule
  SDK maintainers
    & Keep governance and policy control, approve rollout timing, and retain security gates. \\[4pt]
  SDK repositories
    & Keep workflow YAML, permissions, checkout, concurrency, harden-runner, and local config. \\[4pt]
  \texttt{sdk-automations}
    & Provide shared core logic, adapters, schemas, releases, fixtures, and examples. \\[4pt]
  GitHub Actions
    & Execute SDK-owned workflows that call central actions by pinned references. \\[4pt]
  Future GitHub App / Probot
    & Later orchestration path that reuses the same core without a second implementation. \\
  \bottomrule
\end{tabularx}

% ============================================================
\section{Constraints}
% ============================================================

\begin{itemize}
  \item GitHub Actions is the current execution environment; workflows remain in SDK repos.
  \item SDK repos retain security controls: permissions, harden-runner, checkout, concurrency.
  \item Central repo must be org-owned and governed before production dependency.
  \item Repo policy config is protected and must remain local to each SDK.
  \item No privileged tokens or escalated permissions originate from the shared core.
  \item Cross-SDK behavior differences require parity tests before write mode.
\end{itemize}

% ============================================================
\section{Context and Scope}
% ============================================================

The architecture boundary is deliberate: SDK repositories remain the security and governance
boundary, while \texttt{sdk-automations} becomes the shared implementation boundary for reusable
automation logic.

\begin{figure}[h]
  \centering
  \includegraphics[width=\linewidth]{image.png}
  \caption{C4 L1 context view for shared SDK automations (Excalidraw).}
\end{figure}

\textbf{Phase 0 proof:} a \href{https://github.com/darshit2308/hiero-sdk-python/actions/runs/25763950499/job/75671903429}{Python fork canary run} validated the end-to-end call path while keeping
workflow boundaries local and applying expected review labels.

\subsection{Scope}

\textbf{In scope:}
\begin{itemize}
  \item Shared automation decision logic, schemas, fixtures, and adapters.
  \item SDK-owned workflows and protected repo policy configuration.
  \item Incremental adoption by SDKs with rollback by pinned SHA.
\end{itemize}

\begin{notebox}
  \textbf{Out of scope:} removing all SDK workflows, bypassing maintainer approval, moving
  privileged policy into a personal fork, or forcing C++ workflows onto an unproven wrapper.
\end{notebox}

% ============================================================
\section{Solution Strategy}
% ============================================================

Adopt a shared-core with ports-and-adapters architecture aligned to the Microkernel pattern. The
core houses reusable automation decisions, while adapters map runtime inputs from GitHub Actions
and (later) Probot to the same core API. GitHub events trigger the system, but the architecture is
not a hosted microservice first; it is a modular automation kernel with thin delivery adapters.

\subsection{Microkernel pattern fit}

\begin{tabularx}{\linewidth}{@{} L{3.2cm} L{2.2cm} X @{}}
  \toprule
  \textbf{Attribute} & \textbf{Expected fit} & \textbf{Rationale for this system} \\
  \midrule
  Agility
    & High
    & New or changing automations are plug-ins in the core without rewriting runtimes. \\[4pt]
  Testability
    & High
    & Core logic is isolated and verified with fixtures and parity tests. \\[4pt]
  Deployability
    & High
    & Adapters and actions are versioned and pinned independently per SDK. \\[4pt]
  Scalability
    & Medium-Low
    & Workload is bounded by SDK events, not high-throughput service traffic. \\[4pt]
  Simplicity
    & Medium
    & Clear core/adapter boundaries reduce duplication but require adapter discipline. \\
  \bottomrule
\end{tabularx}

% ============================================================
\section{Building Block View}
% ============================================================

\subsection{C4 L2 - Container view (text)}

\begin{tabularx}{\linewidth}{@{} L{4.2cm} X @{}}
  \toprule
  \textbf{Container} & \textbf{Responsibility and interaction} \\
  \midrule
  SDK repositories
    & Own workflows, policy config, triggers, and permissions; call shared actions. \\[4pt]
  \texttt{sdk-automations} repo
    & Hosts core logic, adapters, schemas, actions, fixtures, and releases. \\[4pt]
  GitHub Actions runtime
    & Executes SDK workflows and runs shared actions. \\[4pt]
  Future GitHub App / Probot
    & Alternative runtime that calls the same core API. \\[4pt]
  GitHub API
    & External system for comments, labels, assignments, and PR metadata. \\
  \bottomrule
\end{tabularx}

\subsection{C4 L3 - Component view (text)}

\begin{tabularx}{\linewidth}{@{} L{4.7cm} X @{}}
  \toprule
  \textbf{Component} & \textbf{Role} \\
  \midrule
  \texttt{packages/core}
    & Automation decision kernel, pure logic, no runtime-specific I/O. \\[4pt]
  \texttt{packages/github-action-\allowbreak adapter}
    & Validates inputs, loads config, maps Action runtime to core calls. \\[4pt]
  \texttt{packages/probot-app}
    & Maps webhook events to the same core interface used by Actions. \\[4pt]
  \texttt{actions/review-sync}
    & Thin action entrypoint that calls the Action adapter. \\[4pt]
  \texttt{schemas} and fixtures
    & Shared config schema, validation, and parity test fixtures. \\
  \bottomrule
\end{tabularx}

\subsection{Three-bucket classification}

\begin{tabularx}{\linewidth}{@{} L{3.5cm} L{5cm} X @{}}
  \toprule
  \textbf{Bucket} & \textbf{Examples} & \textbf{Rule} \\
  \midrule
  Shared core candidate
    & Review-sync decisions, assignment eligibility, PR check rules, comment builders,
      label transition decisions, fixtures, parity tests, PR priority scoring
      (issue label -> weight + age factor).
    & Move here when behavior is common or intentionally configurable across SDKs. \\[6pt]
  Repo-local boundary
    & Workflow triggers, permissions, checkout, harden-runner, concurrency,
      artifact download / upload, token choice.
    & Keep local because these are security and operational controls. \\[6pt]
  Repo-specific policy
    & Maintainer team, label names, documentation links, assignment limits,
      support channels.
    & Keep as protected local config validated by the central schema. \\
  \bottomrule
\end{tabularx}

\subsection{Candidate repository shape}

The candidate central repository should expose the following components:

\begin{itemize}
  \item \texttt{packages/core}
  \item \texttt{packages/github-action-adapter}
  \item \texttt{packages/probot-app}
  \item \texttt{actions/review-sync}
  \item \texttt{actions/run-automation}
  \item \texttt{schemas}
  \item \texttt{docs}
  \item \texttt{examples}
  \item \texttt{CI}
\end{itemize}

Personal fork links are proof artifacts only; upstream SDKs should depend only on an approved
org-owned repository with maintainers and release governance.

% ============================================================
\section{Runtime View}
% ============================================================

\subsection{GitHub Actions runtime}

\begin{enumerate}
  \item The SDK workflow starts from its own trigger, permissions, concurrency,
        and maintainer-approved YAML.
  \item The job keeps \texttt{step-security/harden-runner} visible near the start of the job
        and checks out only what the job is allowed to use.
  \item The workflow loads protected repo config from the trusted branch and calls
        the shared action by a reviewed immutable reference.
  \item The Action adapter validates inputs and config, then calls \texttt{packages/core}
        for the automation decision.
  \item The core returns dry-run output or performs allowed writes through the GitHub API,
        depending on the approved workflow mode.
\end{enumerate}

\subsection{Probot runtime}

\begin{enumerate}
  \item A future GitHub App receives the webhook and authenticates as the app installation.
  \item The Probot adapter maps the event into the same core function shape used by Actions.
  \item The core produces the same comments, labels, assignment decisions, and safety checks.
\end{enumerate}

% ============================================================
\section{Deployment View}
% ============================================================

\begin{itemize}
  \item Production callers should pin central actions and third-party actions to
        full-length commit SHAs, using tags only as human-readable release labels.
  \item The central repo must move under the approved official owner before upstream
        SDK production dependency.
  \item Branch protection, CODEOWNERS, release owners, changelog, rollback instructions,
        and incident ownership are part of initial governance, not optional extras.
  \item Repo config is privileged policy and must require protected-branch review because
        changing config changes automation behavior.
\end{itemize}

\textbf{Migration sequence implied by this architecture:}
\begin{itemize}
  \item Phase 0 (complete): Python fork canary proves the caller contract.
  \item Phase 1: Stabilize central repo governance before any upstream dependency.
  \item Phase 2: Python review-sync dry-run pilot upstream (no existing workflow replaced).
  \item Phase 3: Controlled write pilot with named owner and rollback SLA.
  \item Phase 4+: Assignment mapping, then C++ canary, then wider SDK rollout.
\end{itemize}

% ============================================================
\section{Crosscutting Concepts}
% ============================================================

\begin{itemize}
    \item \textbf{Security:} workflow YAML remains local, permissions are explicit, and harden-runner
        stays visible in SDK workflows.
    \item \textbf{Configuration:} per-SDK policy config is validated by shared schemas.
    \item \textbf{Testing:} fixtures and parity tests validate core decisions against existing behavior.
    \item \textbf{Observability:} consistent logs and dry-run outputs for reviewers and maintainers.
    \item \textbf{Versioning:} SDKs pin SHAs and can roll back by reverting a single caller change.
\end{itemize}

% ============================================================
\section{Architectural Decisions}
% ============================================================

\begin{itemize}
  \item \textbf{Microkernel pattern}: the shared core is the automation kernel, adapters are plug-ins.
  \item \textbf{Actions-first, Probot-ready}: use GitHub Actions for early pilots; the
        ports-and-adapters split keeps \texttt{packages/core} stable as runtimes change, so a
        Probot adapter maps webhook events to the same function signatures without core rewrites.
  \item \textbf{Do not centralise workflow YAML}: centralise logic only where duplication creates
        real maintenance cost.
  \item \textbf{Python review-sync first}: start with lower-risk automation already proven in a
        shared-action form.
  \item \textbf{Assignment logic deferred}: wait for a Python/C++ behavior map before extraction.
\end{itemize}

The migration strategy is documented separately and derived from this architecture.

% ============================================================
\section{Quality Requirements}
% ============================================================

\begin{tabularx}{\linewidth}{@{} L{3.5cm} X @{}}
  \toprule
  \textbf{Quality goal} & \textbf{Scenario or test} \\
  \midrule
  Security
    & A contributor-facing workflow can call the shared action without hiding harden-runner,
      expanding token permissions unexpectedly, or checking out untrusted PR code in a
      privileged job. \\[6pt]
  Maintainability
    & A bug fix in shared review-sync or assignment logic is tested once in
      \texttt{packages/core} and released once for SDK repos to adopt deliberately. \\[6pt]
  Evolvability
    & A future Probot / GitHub App adapter can call the same core functions without a
      second implementation of assignment or PR checks. \\[6pt]
  Testability
    & Before write mode, golden fixtures and behavior parity tests prove equivalence with
      existing SDK behavior, not only human dry-run log review. \\[6pt]
  Adoptability
    & Each SDK can adopt one automation at a time, keep local policy, and roll back by
      reverting one caller change or pinning an older approved SHA. \\
  \bottomrule
\end{tabularx}

% ============================================================
\section{Risks and Technical Debt}
% ============================================================

\begin{tabularx}{\linewidth}{@{} L{4.2cm} X @{}}
  \toprule
  \textbf{Risk or debt} & \textbf{Mitigation or note} \\
  \midrule
  Shared core defect impacts multiple SDKs
    & Pin versions, stage rollout, and roll back by SHA; keep parity tests green. \\[4pt]
  Behavior drift between SDKs
    & Maintain behavior maps, per-SDK fixtures, and strict schema validation. \\[4pt]
  Central repo governance delayed
    & Keep local automations active and treat shared core as internal until approved. \\[4pt]
  Security regression from hidden permissions
    & Keep workflow YAML local and require security review for any permission change. \\[4pt]
  C++ wrapper constraints unknown
    & Keep C++ workflows local until a proven wrapper and parity results exist. \\
  \bottomrule
\end{tabularx}

% ============================================================
\section{Glossary}
% ============================================================

\begin{tabularx}{\linewidth}{@{} L{3.5cm} X @{}}
  \toprule
  \textbf{Term} & \textbf{Meaning} \\
  \midrule
  \texttt{sdk-automations}
    & Central repository hosting shared automation logic and adapters. \\[4pt]
  \texttt{packages/core}
    & Automation decision kernel, runtime-agnostic logic. \\[4pt]
  Adapter
    & Thin layer that maps runtime inputs to the core interface. \\[4pt]
  harden-runner
    & StepSecurity action that hardens GitHub Actions runners. \\[4pt]
  dry-run mode
    & Non-writing execution that outputs intended actions for review. \\[4pt]
  parity tests
    & Tests that compare core outputs against existing SDK behavior. \\[4pt]
  review-sync
    & Automation that synchronizes review state, labels, or comments. \\
  \bottomrule
\end{tabularx}

% ============================================================
\section*{References}
% ============================================================

\begin{itemize}
  \item arc42 overview: \url{https://arc42.org/overview}

  \item Architecture documentation guide and C4 framing:\\
        \url{https://www.workingsoftware.dev/software-architecture-documentation-the-ultimate-guide/}

  \item Richards, Mark. \textit{Software Architecture Patterns}. O'Reilly, 2015.

  \item GitHub Actions secure use reference:\\
        \url{https://docs.github.com/en/actions/reference/security/secure-use}

  % \item GitHub \texttt{pull\_request\_target} warning:\\
  %       \url{https://docs.github.com/en/actions/writing-workflows/choosing-when-your-workflow-runs/events-that-trigger-workflows}

  \item StepSecurity Harden-Runner docs: \url{https://docs.stepsecurity.io/harden-runner}
\end{itemize}

% ============================================================
\end{document}
% ============================================================